% Ad Atticum 13.12 (ad-atticum-13-12)
% Source: https://raw.githubusercontent.com/PerseusDL/canonical-latinLit/master/data/phi0474/phi057/phi0474.phi057.perseus-lat1.xml
% Fetched by scripts/fetch_latin.py

% Dateline (Perseus): Scr. in Arpinati vh i K. Quint. a. 709 (45).

\ciceroLetterOpener{CICERO ATTICO salutem}

\ciceroSection{1} valde me memorderunt epistulae tuae de Attica nostra; eaedem tamen sanaverunt. quod enim te ipse consolabare eisdem litteris, id mihi erat satis firmum ad leniendam aegritudinem.

\ciceroSection{2} Ligarianam praeclare vendidisti. posthac quicquid scripsero tibi praeconium deferam.

\ciceroSection{3} quod ad me de Varrone scribis, scis me antea orationes aut aliquid id genus solitum scribere ut Varronem nusquam possem intexere. postea autem quam haec coepi φιλολογώτερα , iam Varro mihi denuntiaverat magnam sane et gravem προσφώνησιν . biennium praeteriit cum ille Καλλιπίδησ adsiduo cursu cubitum nullum processerit, ego autem me parabam ad id quod ille mihi misisset ut αὐτῷ τῷ μέτρῳ καὶ , si modo potuissem; nam hoc etiam Hesiodus ascribit, αἴ κε δύνηαι . nunc illam περὶ Τελῶν σύνταξιν sane mihi probatam Bruto, ut tibi placuit, despondimus, idque tu eum non nolle mihi scripsisti. ergo illam Ἀκαδημικήν , in qua homines nobiles illi quidem sed nullo modo philologi nimis acute loquuntur, ad Varronem transferamus. etenim sunt Antiochia quae iste valde probat. Catulo et Lucullo alibi reponemus, ita tamen si tu hoc probas; deque eo mihi rescribas velim.

\ciceroSection{4} de Brinniana auctione accepi a Vestorio litteras. ait sine ulla controversia rem ad me esse conlatam. Romae videlicet aut in Tusculano me fore putaverunt a. d. viii Kal. Quint. dices igitur vel amico tuo S. Vettio coheredi meo vel Labeoni nostro paulum proferant auctionem; me circiter Nonas in Tusculano fore. cum Pisone Erotem habes. de Scapulanis hortis toto pectore cogitemus. dies adest.
